% =====================================================================
\section{Multimodal Model Architecture}
\label{sec:model}
% =====================================================================

This section specifies the complete recognition model in enough detail to be
implemented directly. It combines the encoders of \cref{sec:visual,sec:audio}
with a fusion module and a CTC output layer. \Cref{tab:model-summary} gives the
recommended dimensions \rec; all are starting values to be tuned.

\begin{table}[H]
\centering
\caption{Summary of the proposed model (recommended starting configuration).}
\label{tab:model-summary}
\small
\begin{tabularx}{\textwidth}{@{}p{3.2cm}Xl@{}}
\toprule
\textbf{Component} & \textbf{Configuration} & \textbf{Output} \\
\midrule
Visual encoder & 3 Conv3d blocks (32/64/96 ch.), Linear 512, 2-layer BiLSTM (256/dir.) & \shape{B,T,512} \\
Audio encoder & 2 strided Conv1d (256 ch.), 2-layer BiLSTM (256/dir.) & \shape{B,T,512} \\
Projections & Linear $512{\to}256$ per stream, LayerNorm & \shape{B,T,256} each \\
Cross-modal attention & Two MHA blocks (A$\to$V, V$\to$A), 4 heads, $d=256$ & \shape{B,T,256} each \\
Gated fusion & Per-frame scalar gate $g_t$ & \shape{B,T,256} \\
Sequence model & 1-layer BiLSTM (128/dir.) & \shape{B,T,256} \\
CTC output layer & Linear $256{\to}28$, log-softmax & \shape{T,B,28} \\
\bottomrule
\end{tabularx}
\end{table}

\subsection{Visual encoder}
\textbf{Purpose.} Convert a mouth-crop sequence into one feature vector per
video frame that describes the lip configuration and its motion.
\textbf{Input.} \shape{B,1,T,96,96} normalised crops and a frame mask.
\textbf{Operations.} 3D convolutions apply kernels that extend over three
consecutive frames as well as over space. A single kernel can therefore respond
to a pattern such as ``lips closing'', which no single frame shows; this is how
the encoder captures the motion that distinguishes otherwise similar mouth
shapes~\cite{lipsyncnet2024,li2023}. Spatial pooling progressively removes
position detail, and the BiLSTM integrates evidence across the whole
utterance, using both past and future frames.
\textbf{Output.} $\mathbf{H}_v \in \mathbb{R}^{B\times T\times 512}$.
\textbf{Integration.} Passed to the fusion module; also used alone by the
visual-only baseline.
\textbf{Testing.} Shape tests (\cref{sec:visual}); a visual-only model must
overfit a batch of 8 utterances to near-zero CTC loss.

\subsection{Audio encoder}
\textbf{Purpose.} Convert acoustic features into one vector per video frame
describing the sounds being produced.
\textbf{Input.} \shape{B,T_a,80} standardised log-Mel features and lengths.
\textbf{Operations.} Two strided 1-D convolutions reduce 100\,fps features to
25\,fps (\cref{sec:audio-encoder}); a two-layer BiLSTM adds utterance-level
context.
\textbf{Output.} $\mathbf{H}_a \in \mathbb{R}^{B\times T\times 512}$.
\textbf{Integration.} Passed to fusion; used alone by the audio-only baseline.
\textbf{Testing.} Shape and length tests (\cref{sec:audio}); overfitting test as
for the visual encoder.

\subsection{Feature alignment and fusion}
\label{sec:fusion}

\paragraph{Temporal alignment.} The two streams have different native rates:
25\,fps for video and 100\,fps for audio features. Rather than resampling
inside the fusion module, the audio encoder subsamples by exactly four, so both
encoders emit sequences of the same length $T$ on the same 40\,ms time grid.
Frame $t$ of $\mathbf{H}_a$ and frame $t$ of $\mathbf{H}_v$ therefore describe the
same instant. If the two lengths differ by one or two frames (rounding at clip
boundaries), the longer sequence is truncated to the shorter. Cross-attention
itself does not require equal lengths, but the per-frame gate below does.

\paragraph{Projection.} Each stream is projected to a shared width $d=256$:
\begin{equation}
\mathbf{A} = \operatorname{LN}(\mathbf{H}_a\mathbf{W}_a), \qquad
\mathbf{V} = \operatorname{LN}(\mathbf{H}_v\mathbf{W}_v),
\qquad \mathbf{A},\mathbf{V}\in\mathbb{R}^{T\times d},
\end{equation}
where LN is layer normalisation (batch index omitted).

\paragraph{Cross-modal attention.} Each stream queries the other. For audio
attending to video, the \emph{queries} come from the audio sequence and the
\emph{keys} and \emph{values} from the video sequence:
\begin{gather}
\mathbf{Q} = \mathbf{A}\mathbf{W}_Q,\quad \mathbf{K} = \mathbf{V}\mathbf{W}_K,\quad \mathbf{U} = \mathbf{V}\mathbf{W}_U, \\
\boldsymbol{\alpha} = \operatorname{softmax}\!\Big(\frac{\mathbf{Q}\mathbf{K}^{\top}}{\sqrt{d_k}} + \mathbf{M}\Big),
\qquad
\operatorname{Attn}(\mathbf{A},\mathbf{V}) = \boldsymbol{\alpha}\,\mathbf{U}, \\
\tilde{\mathbf{A}} = \operatorname{LN}\big(\mathbf{A} + \operatorname{MHA}(\mathbf{A},\mathbf{V})\big),
\qquad
\tilde{\mathbf{V}} = \operatorname{LN}\big(\mathbf{V} + \operatorname{MHA}(\mathbf{V},\mathbf{A})\big).
\end{gather}
Row $t$ of the attention weights $\boldsymbol{\alpha}\in\mathbb{R}^{T\times T}$
says how strongly audio frame $t$ draws on each video frame; $\mathbf{M}$ sets
padded positions to $-\infty$. MHA denotes the multi-head version with four heads
of width $d_k=64$, implemented with \texttt{nn.MultiheadAttention}. The audio
stream can thereby borrow visual evidence for an ambiguous sound, and the
visual stream can borrow acoustic evidence for an ambiguous lip shape---the
mechanism proposed for disambiguating homophenes~\cite{li2023}.

\paragraph{Gated fusion.} A learned gate decides, frame by frame, how much each
enriched stream contributes:
\begin{equation}
g_t = \sigma\big(\mathbf{w}_g^{\top}[\tilde{\mathbf{a}}_t ; \tilde{\mathbf{v}}_t] + b_g\big) \in (0,1),
\qquad
\mathbf{f}_t = g_t\,\tilde{\mathbf{a}}_t + (1-g_t)\,\tilde{\mathbf{v}}_t ,
\label{eq:gate}
\end{equation}
where $[\cdot\,;\cdot]$ denotes concatenation and $\sigma$ the logistic function.
A small $g_t$ means that the fused representation relies mainly on the
video-anchored stream. The fused sequence $\mathbf{F}=(\mathbf{f}_1,\dots,\mathbf{f}_T)$
passes through a one-layer BiLSTM and the CTC output layer.

\paragraph{How noise adaptation is encouraged.} The proposal states that fusion
should give more weight to the lips when the audio is noisy \spec. An attention
module does not detect noise by itself; it learns whatever weighting reduces
the training loss on the data it sees. Three mechanisms are possible:
\begin{enumerate}[label=(\alph*)]
  \item \emph{explicit SNR conditioning}, feeding the true SNR to the gate;
  \item \emph{a learned noise-estimation signal}, an auxiliary head trained to
        predict the SNR from the audio encoder's output;
  \item \emph{implicit learning}, training the attention and gate on a mixture
        of clean and noisy examples so that they learn to down-weight
        unreliable audio.
\end{enumerate}
This plan adopts \textbf{(c)} as the primary mechanism \rec. Option~(a) is not
usable in deployment, because the true SNR of live audio is unknown. Training on
all five SNR conditions (\cref{sec:noise}) exposes the gate to examples in which
the audio is uninformative, and modality dropout (\cref{sec:modality-dropout})
adds examples in which it is absent. Option~(b) is kept as an experiment
\tbd: an auxiliary linear head on the time-averaged audio encoding regresses the
known training SNR (with ``clean'' mapped to 30\,dB), its prediction is appended
to the gate input in \cref{eq:gate}, and its loss is added to the CTC loss with a
small weight such as 0.1.

\paragraph{How noise adaptation is evaluated.} Whether the gate really behaves
as intended is an empirical question. During testing the evaluation script
records the mean gate value $\bar{g}$ for every utterance and reports its
distribution at each SNR. If the mechanism works, $\bar{g}$ should decrease as
SNR falls, and the fused model's advantage over the audio-only baseline should
grow at 0\,dB and $-5$\,dB. Because $\tilde{\mathbf{a}}_t$ already contains some visual
information through cross-attention, $g_t$ is only an approximate measure of
reliance on audio; the modality ablations of \cref{sec:evaluation}
(removing one stream at test time) provide the more direct evidence.

\subsection{Modality dropout}
\label{sec:modality-dropout}
\textbf{Purpose.} Audio is usually far more informative than video, so a fused
model can learn to ignore the visual stream entirely. Modality dropout forces it
to make use of both streams by occasionally removing one during
training~\cite{afouras2022}.

\textbf{Implementation.} For each training utterance a value $u \sim
\mathcal{U}(0,1)$ is drawn:
\begin{equation}
\text{stream used} =
\begin{cases}
\text{video only (audio dropped)} & u < p_a,\\
\text{audio only (video dropped)} & p_a \le u < p_a + p_v,\\
\text{both} & \text{otherwise}.
\end{cases}
\end{equation}
Both streams are never dropped together. A dropped stream's encoder output is
replaced by a learned ``missing modality'' vector $\mathbf{e}_{\text{miss}}\in\mathbb{R}^{256}$
repeated over all frames, rather than by an attention mask. Masking every key of
a sequence would make the softmax in the cross-attention undefined (all logits
$-\infty$) and produce NaN values; a learned vector avoids this and lets the gate
learn to recognise the missing stream.

\textbf{Rate selection.} Starting values are $p_a=0.25$ and $p_v=0.15$ \rec;
audio is dropped more often because the risk being addressed is under-use of
video. The rates are tuned on validation WER averaged over all five SNR
conditions, using the grid $p_a,p_v \in \{0, 0.1, 0.25, 0.4\}$ with $p_a+p_v\le0.5$
if compute permits. The setting $p_a=p_v=0$ is the ``no modality dropout''
ablation.

\textbf{Inference behaviour.} Modality dropout is disabled at inference. Both
streams are used whenever available. If the face is not found in more than half
of a window's frames, the visual encoding is replaced by $\mathbf{e}_{\text{miss}}$---
the same representation the model saw during training---so the system degrades
to audio-only recognition instead of being fed meaningless crops.

\subsection{Sequence prediction and decoding}
\label{sec:ctc}

\paragraph{Why CTC.} The training data provide a transcript for each utterance
but not the frame at which each character is spoken. Connectionist Temporal
Classification~\cite{graves2006} removes the need for such frame-level labels
by summing over every possible alignment between the $T$ output frames and the
$L$ target characters, as used in the CTC-based lipreading and AVSR systems
cited in the proposal~\cite{kanavnet2026,afouras2022}.

\paragraph{Vocabulary, blank and collapsing.} The output layer predicts a
distribution over the 28 symbols of \cref{eq:vocab} at every frame. CTC
introduces the \emph{blank} symbol (index~0) to mean ``no new character''. A
frame-level path $\boldsymbol{\pi}=(\pi_1,\dots,\pi_T)$ is mapped to text by a
function $\mathcal{B}$ that first merges consecutive repeated symbols and then
removes blanks, for example
$\mathcal{B}(\texttt{bb\_ii\_nn\_\_}) = \texttt{bin}$ and
$\mathcal{B}(\texttt{gr\_e\_e\_n}) = \texttt{green}$.
A doubled letter therefore requires a blank between its two copies.

\paragraph{Loss.} With $y_k^t$ the predicted probability of symbol $k$ at frame
$t$, the probability of transcript $\mathbf{z}$ given input $\mathbf{x}$ and the
CTC loss are
\begin{equation}
p(\mathbf{z}\mid\mathbf{x}) = \sum_{\boldsymbol{\pi}\in\mathcal{B}^{-1}(\mathbf{z})}\;\prod_{t=1}^{T} y^{t}_{\pi_t},
\qquad
\mathcal{L}_{\text{CTC}} = -\log p(\mathbf{z}\mid\mathbf{x}).
\label{eq:ctc}
\end{equation}
The sum is computed efficiently by the forward--backward algorithm. The
implementation uses PyTorch's built-in loss,
\texttt{nn.CTCLoss(blank=0, zero\_infinity=True)}, called with
log-probabilities of shape \shape{T,B,28}, the concatenated target indices, the
per-utterance encoder lengths $T^{(b)}$ and target lengths $L^{(b)}$. The
\texttt{zero\_infinity} flag prevents an occasional impossible alignment from
producing an infinite loss that would corrupt training.

\paragraph{Conditions for CTC to work.} An alignment exists only if
\begin{equation}
T^{(b)} \;\ge\; L^{(b)} + R^{(b)},
\label{eq:ctc-cond}
\end{equation}
where $R^{(b)}$ is the number of adjacent repeated characters in the target (each
needing an extra blank). GRID transcripts contain at most about 35 characters
including spaces, while a 3\,s clip gives $T=75$ encoder frames, so the condition
holds with a wide margin. It would fail if the audio encoder subsampled more
aggressively or if live windows were made very short; the data loader therefore
asserts \cref{eq:ctc-cond} for every item, and the window length of the live
demo is chosen with it in mind. Encoder lengths must also be computed from the
true (unpadded) input lengths using the convolution length formula of
\cref{sec:audio-encoder}; passing padded lengths is a common silent error.

\paragraph{Greedy decoding.} Taking the most probable symbol at each frame and
applying $\mathcal{B}$ is fast and is used for validation during training.

\paragraph{Beam-search decoding.} The proposal specifies beam search for final
decoding \spec. CTC prefix beam search keeps the $W$ most probable \emph{text
prefixes}, tracking for each prefix separately the probability of paths ending
in a blank and in a non-blank symbol, so that paths that collapse to the same
text are merged. A beam width of $W=10$ is the starting value \rec. With only
28 symbols the algorithm is implemented directly in about 60 lines of Python,
which keeps it transparent and testable.

\paragraph{Language-model integration.} An external language model could be added
by rescoring beam prefixes. On GRID this needs care: the fixed six-slot grammar
means that a grammar-aware language model would correct many errors by
construction, inflating accuracy in a way that would not carry over to real
captioning. Headline results will therefore be reported \emph{without} a
language model \rec. A grammar-constrained decoder may be reported separately,
clearly labelled, as an upper reference \tbd.

\subsection{Forward pass and decoding algorithms}

\begin{algorithm}[H]
\caption{LipScribe forward pass (training and inference).}
\label{alg:forward}
\begin{algorithmic}[1]
\Require mouth crops $X_v$ \shape{B,1,T,96,96}, log-Mel features $X_a$ \shape{B,T_a,80}, lengths $T_a^{(b)}$, flag \textit{training}
\Ensure log-probabilities $\mathbf{P}$ \shape{T,B,28}, encoder lengths $T^{(b)}$, gate values $g$
\State $\mathbf{H}_v \gets \textsc{VisualEncoder}(X_v)$ \Comment{\shape{B,T,512}}
\State $\mathbf{H}_a \gets \textsc{AudioEncoder}(X_a, T_a^{(b)})$; \quad $T^{(b)} \gets \lceil T_a^{(b)}/4\rceil$ \Comment{\shape{B,T,512}}
\State $\mathbf{A} \gets \operatorname{LN}(\mathbf{H}_a\mathbf{W}_a)$; \quad $\mathbf{V} \gets \operatorname{LN}(\mathbf{H}_v\mathbf{W}_v)$
\If{\textit{training}} apply modality dropout: replace $\mathbf{A}$ or $\mathbf{V}$ by $\mathbf{e}_{\text{miss}}$ per sample
\ElsIf{face missing in window} $\mathbf{V} \gets \mathbf{e}_{\text{miss}}$
\EndIf
\State $\tilde{\mathbf{A}} \gets \operatorname{LN}(\mathbf{A} + \operatorname{MHA}(\mathbf{A},\mathbf{V},\text{mask}))$; \quad
       $\tilde{\mathbf{V}} \gets \operatorname{LN}(\mathbf{V} + \operatorname{MHA}(\mathbf{V},\mathbf{A},\text{mask}))$
\State $g \gets \sigma(\mathbf{w}_g^\top[\tilde{\mathbf{A}};\tilde{\mathbf{V}}] + b_g)$; \quad $\mathbf{F} \gets g\odot\tilde{\mathbf{A}} + (1-g)\odot\tilde{\mathbf{V}}$
\State $\mathbf{Z} \gets \textsc{BiLSTM}(\mathbf{F}, T^{(b)})$; \quad $\mathbf{P} \gets \operatorname{log\_softmax}(\mathbf{Z}\mathbf{W}_o + \mathbf{b}_o)$, transposed to time-first
\State \Return $\mathbf{P}, T^{(b)}, g$
\end{algorithmic}
\end{algorithm}

\begin{algorithm}[H]
\caption{CTC prefix beam search for one utterance (probabilities in log space in practice).}
\label{alg:beam}
\begin{algorithmic}[1]
\Require frame posteriors $y^t_k$ for $t=1..T$, beam width $W$, blank index $0$
\Ensure most probable transcript
\State $\mathcal{P} \gets \{\varepsilon\}$; \quad $p_b(\varepsilon) \gets 1$, $p_{nb}(\varepsilon) \gets 0$
\For{$t = 1$ \textbf{to} $T$}
  \State initialise $p'_b, p'_{nb}$ to 0 for all prefixes
  \For{each prefix $\ell \in \mathcal{P}$}
    \State $p'_b(\ell) \mathrel{+}= y^t_0\,(p_b(\ell)+p_{nb}(\ell))$ \Comment{emit blank}
    \For{each symbol $k \ne 0$}
      \State $\ell^{+} \gets \ell \Vert k$
      \If{$k$ = last symbol of $\ell$}
        \State $p'_{nb}(\ell^{+}) \mathrel{+}= y^t_k\,p_b(\ell)$ \Comment{repeat needs a blank in between}
        \State $p'_{nb}(\ell) \mathrel{+}= y^t_k\,p_{nb}(\ell)$ \Comment{repeat collapses}
      \Else
        \State $p'_{nb}(\ell^{+}) \mathrel{+}= y^t_k\,(p_b(\ell)+p_{nb}(\ell))$
      \EndIf
    \EndFor
  \EndFor
  \State $\mathcal{P} \gets$ the $W$ prefixes with largest $p'_b + p'_{nb}$; \quad $p_b, p_{nb} \gets p'_b, p'_{nb}$
\EndFor
\State \Return $\arg\max_{\ell\in\mathcal{P}} \big(p_b(\ell)+p_{nb}(\ell)\big)$ mapped to characters
\end{algorithmic}
\end{algorithm}
